\section*{Begrebsforklaring}
\label{sec:begreber}
\addcontentsline{toc}{section}{Begrebsforklaring}

The table below centralises core notation and concepts used across Part I and Part II.

\begin{table}[H]
\centering
\footnotesize
\caption{Core notation and modelling concepts used in this report}
\label{tab:begreber_notation}
\rowcolors{3}{black!4}{white}
\begin{tabularx}{\textwidth}{>{\raggedright\arraybackslash}p{0.24\textwidth} X}
\toprule
Term & Explanation \\
\midrule
Axx & Activity ID used in the interview-based routine model (Part I). \\
Exx & Event ID used in the observed event log (Part I). \\
Dxx & Delta ID used in the comparison sheet between stated and observed routine (Part I). \\
Mx & Method ID used for toolkit methods in the Part II notebook workflow. \\
DCR & Dynamic Condition Response process modelling notation. \\
Interview-based DCR & DCR representation built from interview evidence (intended routine structure). \\
Observed DCR & DCR evidence grounded in observed/logged events. \\
Say-do gap & Difference between how work is described and how it is actually performed. \\
Delta sheet & Structured comparison table capturing mismatches between stated and observed process traces. \\
Manual knowledge acquisition & Knowledge elicitation through qualitative methods such as interviews and observation. \\
Knowledge-as-stated & Knowledge captured from participant description/interview account. \\
Knowledge-as-discovered & Knowledge revealed through observation and comparison of empirical traces. \\
Work-as-imagined & Model of how the process is expected or intended to run. \\
Work-as-done & What is actually performed in practice. \\
Work-as-recorded & What is represented in available data/log traces. \\
\bottomrule
\end{tabularx}
\rowcolors{3}{}{}
\end{table}

\begin{table}[H]
\centering
\footnotesize
\caption{Part II data-quality and analysis concepts}
\label{tab:begreber_part2}
\rowcolors{3}{black!4}{white}
\begin{tabularx}{\textwidth}{>{\raggedright\arraybackslash}p{0.24\textwidth} X}
\toprule
Term & Explanation \\
\midrule
RQ & Research question used to contextualise data quality assessment decisions. \\
RQ1 & Workflow efficiency and bottlenecks, including turnaround-time interpretation. \\
RQ2 & Timeliness of critical finding communication for patient-safety interpretation. \\
RQ3 & AI-based triage use-case and data suitability for modelling. \\
DQ & Data quality. \\
Five Facets framework & Data quality framing across Data, Source, System, Task, and Human facets. \\
Data facet & Quality properties of captured data values and structure. \\
Source facet & How source processes and contexts shape data generation. \\
System facet & Effects of technical systems/interfaces on what is captured and how. \\
Task facet & Alignment between data capture and clinical/operational task requirements. \\
Human facet & Human behaviour, workload, and documentation practice influences on data quality. \\
Completeness & Presence/absence of required values and fields. \\
Representativity & Degree to which data coverage reflects relevant populations, contexts, and workload patterns. \\
Conformance check & Verification that logged traces follow expected process or timing rules. \\
Timeliness & Whether event timing supports valid performance or safety interpretation. \\
TAT & Turnaround time; elapsed time between selected workflow milestones. \\
Batch logging & Retrospective or clustered timestamp entry that can distort temporal metrics. \\
Invisible work & Clinically relevant activities present in practice/workflow but absent from event logs. \\
\bottomrule
\end{tabularx}
\rowcolors{3}{}{}
\end{table}
